\documentclass[11pt]{article}
\usepackage[T1]{fontenc}
\usepackage[utf8]{inputenc}
\usepackage{lmodern,amsmath,amssymb,microtype,xcolor}
\usepackage[margin=.9in]{geometry}
\usepackage[colorlinks=true,urlcolor=blue!60!black]{hyperref}
\setlength{\parindent}{0pt}
\setlength{\parskip}{.5em}
\setlength{\emergencystretch}{2em}
\author{Wentao Li\\\small Prepared with assistance from OpenAI Codex}
\date{November 22, 2026}
\title{Arithmetic progressions in the image of an infinite clique}
\hypersetup{pdfauthor={Wentao Li},pdftitle={Arithmetic progressions in the image of an infinite clique}}
\begin{document}
\maketitle
Filipów, Kowitz and Kwela ask whether the van der Waerden ideal is Katětov below the Ramsey ideal and the difference ideal. Both comparisons fail in ZFC. The obstruction is stronger than bounded progression length: an infinite clique can always be thinned until its edge labels contain no genuine three-term arithmetic progression. \href{https://arxiv.org/abs/2307.06881}{2023 preprint, Question 1.2} · \href{https://doi.org/10.1017/jsl.2024.8}{Unified paper, Question 7.8 and Question 16.6}

\textbf{Theorem.} Let \(f:[\omega]^2\to\mathbb R\), and let \(\mathcal L\) be a finite list of triples \((a,b,c)\) of nonzero real numbers with \(a+b+c=0\). There is an infinite \(B\subseteq\omega\) such that no three distinct values in \(f[[B]^2]\) obey any of the specified equations \(au+bv+cw=0\). The proof works for arbitrary edge labels. In particular it applies to \(1,-2,1\).

\textbf{Proof.} The canonical Ramsey theorem for pairs gives an infinite set on which equality of labels has one of four forms: all equal, equality of first coordinate, equality of last coordinate, or equality of the entire ordered pair (write every edge in increasing order). In the constant case there is nothing to prove. In either one-coordinate case the sequence of distinct labels has an infinite subsequence avoiding all the equations. At each step, two old values, their positions and the coefficient triple determine at most one forbidden new value. There are only finitely many choices. Use the corresponding subsequence of vertices.

In the injective case thin the vertices by the ordinary infinite Ramsey theorem finitely many times. For each ordered pattern of three distinct edges on at most six vertices and each assignment of the three coefficients, make the truth value of the corresponding equation constant on all copies of the pattern. There are only finitely many colorings and each is a two-coloring of subsets of a fixed finite size. Suppose some homogeneous pattern satisfies an equation.

If its underlying graph has a vertex of degree one, embed the pattern with enough gaps in the infinite ordered set to replace that vertex while leaving all others fixed. Exactly one labeled edge changes. Its coefficient is nonzero and injectivity changes its value, contradicting the equation. A simple graph on its nonisolated vertices with three distinct edges and minimum degree at least two is a triangle. Thus for some permutation of the coefficients, again called \(a,b,c\), one must have \[
 a f(i,j)+b f(i,k)+c f(j,k)=0 \quad\text{for every }i<j<k
\] in the thinned set. Fix five vertices \(i<j<m<k<l\). For \(v\in\{i,j,m\}\) write \(D_v=f(v,k)-f(v,l)\ne0\). Subtraction gives \(bD_v+cD_w=0\) for every \(v<w\) from the first three vertices. The first two such equations force \(D_j=D_m\), and the third gives \(b+c=0\). Then \(a=0\), a contradiction. All homogeneous patterns therefore avoid the equations.

\textbf{Ideal and topological consequences.} Write \(\mathcal W\) for the sets of natural numbers with a finite bound on their arithmetic progression lengths, \(\mathcal R\) for graphs containing no infinite clique, and \(\mathcal D\) for sets containing no \(\Delta(B)=\{j-i:i<j,\ i,j\in B\}\) with \(B\) infinite. The convention \(\mathcal I\le_K\mathcal J\) requires a map on the domain of \(\mathcal J\) whose inverse image takes each member of \(\mathcal I\) into \(\mathcal J\). For a candidate map from edges to naturals, its set of labels on the thinned clique is in \(\mathcal W\) but the preimage contains the clique. This rules out \(\mathcal W\le_K\mathcal R\). Apply the same argument to \(f(i,j)=g(j-i)\) for any candidate \(g:\omega\to\omega\). The preimage of the progression-free image now contains \(\Delta(B)\), ruling out \(\mathcal W\le_K\mathcal D\).

Under CH the characterization of Filipów, Kowitz and Kwela (Theorem 16.1 and Corollary 16.7) transfers these obstructions to spaces. In particular there exists a compact separable Hausdorff Ramsey space that is not a van der Waerden space. Here their Ramsey function is \(P^-\) and the van der Waerden ideal is \(F_\sigma\), hence \(P^+\) for the theorem. Every Ramsey space is differentially compact by pulling back a difference-indexed sequence to pairs. Thus the same example separates differential compactness from van der Waerden compactness. These applications answer the Ramsey inclusion question by a consistency counterexample. We have no unconditional construction of that separating space here. \href{https://doi.org/10.1017/jsl.2024.8}{Filipów, Kowitz and Kwela, §16}

\textbf{Scope and priority.} The canonization of equality types is the classical Erdős--Rado theorem. The argument here canonizes the linear relations too and excludes the remaining triangle configuration. The Hindman-ideal comparison in the same source asks about all finite sums and does not follow from the pair argument. Provisional significance \textbf{59/100}, conditional on correctness and novelty. The two ideal comparisons use one theorem. Independent proof and priority review are pending.
\end{document}
